\chapter{Resources}

\section*{Further reading} \phantomsection \addcontentsline{toc}{section}{Further reading}

\subsection*{The flawed human brain}

\subsection*{Thinking, Fast and Slow, Daniel Kahneman, 2011, Penguin}

Excellent book on \textbf{cognitive bias.} A must read.

\subsection*{Beyond Greed and Fear, Hersh Shefrin, 2007, OUP USA}

Relatively short book on \textbf{cognitive bias} in finance specifically. Worth reading if you don't have time for Thinking, Fast and Slow.

\subsection*{The Education of a Speculator, Victor Niederhoffer, 1998, Wiley}

Fascinating and esoteric book by a famous partly systematic, negative skew, hedge fund manager. Part autobiography, part book on the philosophy of trading. Optional reading.

\subsection*{Systematic trading rules}

\subsection*{More Money than God, Sebastian Mallaby, 2011, Penguin}

A history of hedge funds, but also a very readable guide to different strategies. Compulsory reading.

\subsection*{Expected Returns, Antti Ilmanen, 2011, Wiley Finance}

Comprehensive guide to the sources of returns and risk premia. Read after Mallaby, if you can cope with the sometimes technical treatment.

\subsection*{When Genius Failed: The Rise and Fall of Long-Term Capital Management, Roger Lowenstein, 2001, Random House}

A \textbf{negative skew} disaster and a real-life example of the disadvantage of leveraging up too much on low volatility positions. Probably the most useful market history in this list and should be compulsory reading for anyone contemplating negative skew trading.

\subsection*{Rogue Trader, Nick Leeson, 1999, Sphere}

Very famous example of when negative skew goes wrong; amongst other things Nick lost much of his money selling option straddles. Optional reading.

\subsection*{The Greatest Trade Ever, Gregory Zuckerman, 2010, Penguin}

A great story about a positive skew trade that worked: John Paulson's bearish bet on mortgage backed securities. Optional reading.

\subsection*{The Black Swan, Nassim Taleb, 2008, Penguin}

A book about unknown unknowns. Taleb's usual mixture of unique philosophy and market folklore. Interesting, but optional reading.

\subsection*{Trading rule fitting}

\subsection*{Fooled by Randomness, Nassim Taleb, 2001, Penguin}

A very interesting book on uncertainty in general. Compulsory for anyone who thinks back-testing is worthwhile.

\subsection*{Trading rules and forecasts}

\subsection*{Trading Systems and Methods, 5th Edition, Perry J. Kaufman, 2013, John Wiley \& Sons}

The bible of trading strategies. Great resource for trading rule ideas if you need them.

\subsection*{Technical Analysis, Jack Schwager, 1995, John Wiley \& Sons}

This, and the next book, are the best in a large crop of similar books. Compulsory if you want ideas for new technical trading rules.

\subsection*{Fundamental Analysis, Jack Schwager, 1997, John Wiley \& Sons}

See above. Compulsory if you want to trade fundamentals.

\subsection*{Hedge Fund Market Wizards, Jack Schwager, 2012, John Wiley \& Sons}

Interviews with many successful hedge fund managers. There are many useful nuggets of information in here. The chapter on Michael Platt is the most relevant to systematic traders. Optional reading.

\subsection*{Volatility targeting}

\subsection*{Fortune's Formula, William Poundstone, 2005, Hill \& Wang}

Highly readable story about the \textbf{Kelly criterion}, and the history of finance in general. Compulsory reading.

\audiencebox{Semi-automatic trader}{\textit{How to Win at Financial Spread Betting}, Charles Vincent, 2002, FT Prentice Hall. \par This is one of several books which provide a good introduction to the mechanics of spread betting in its first half. I don't advise using the strategies in the second half, but then I would say that!}

\audiencebox{Asset allocating investor}{\textit{FT Guide to Exchange Traded Funds and Index Funds}, David Stevenson, 2012, Financial Times. \par UK book about ETFs.
\par
\medskip
\textit{The ETF Book}, Richard Ferri, 2009, Wiley. \par US book about ETFs.}

\audiencebox{Staunch systems trader}{\textit{Trading Commodities and Financial Futures}, George Kleinman, 2013, Prentice Hall. \par One of many books that give a good introduction to futures trading. There are also some ideas for trading rules. A fine alternative would be any of the Schwager books mentioned above.}

\section*{Sources of free data} \phantomsection \addcontentsline{toc}{section}{Sources of free data}

For professional investors sourcing data is usually straightforward, but it can be harder for amateurs, especially those not wishing to pay. Sources of historical data that I've used are listed below. Website links are correct at the time of writing, but are subject to change.

Many brokers and exchanges provide access to live prices, usually with a 15 minute delay. This delay is not problematic unless you are trading on a high frequency basis or using execution algorithms.

\url{www.quandl.com}

Source of numerous kinds of price and fundamental data, both free and subscription. Has a powerful API to automate data collection, available for multiple languages. For my own trading I currently use a combination of data from my broker and Quandl.

\url{finance.yahoo.com}

Source of equity and index data. Has an API to automate data collection.

\url{www.bloomberg.com/markets/chart/data/1D/AAPL:US}

One of the main providers of costly data for professional investors. As the link shows you can get a certain amount of free historic data from Bloomberg, if you know the reference code.

\url{www.eoddata.com}

Free/subscription source of equity data.

\url{www.ivolatility.com}

Source of option prices.

\url{www.oanda.com/currency/historical-rates}

Source of historic FX rates.

\url{stats.oecd.org}

Official source of macroeconomic data.

\url{mba.tuck.dartmouth.edu/pages/faculty/ken.french/data_library.html}

Academic source of equity value data.

\section*{Brokers and platforms} \phantomsection \addcontentsline{toc}{section}{Brokers and platforms}

If you are not running a fully automated trading system, I hesitate to recommend a specific broker. A key selection criteria is price. Estimate the expected annual total fee from how often you're trading, in what size, and with what total account value. Then use this for comparison purposes.

Equally important is whether you can trade the products you want and with leverage if needed. Customer service should be a factor, but it is hard to evaluate in advance. Generally all other features are of secondary importance; although brokers seem to love jazzing up their platforms a flash website won't make you any more money.

For those running an automated system there is only one broker accessible to amateur clients that I can recommend at the time of writing and that is Interactive Brokers (which I use myself). This is because it offers a flexible API which can be used to fully automate your trading and although it is not easy to get it working there is some help on my website: \url{www.systematictrading.org}

There are online platforms that offer to implement your automated strategies for you, as distinct from software that runs on your desktop, which I discuss below. I haven't tested any of these but they look intriguing. Just be very careful, as you are handing over your money to someone else's software!

There are also now ‘social trading' online platforms that allow you to put your money into other amateurs' systematic trading strategies. I would strongly advise against this.

\section*{Automation and coding} \phantomsection \addcontentsline{toc}{section}{Automation and coding}

It is perfectly possible to run non-automated strategies using only simple spreadsheets. There are three avenues for implementing fully automated trading strategies. Firstly it is possible to use Excel to perform automated trading with Interactive Brokers, though this requires some knowledge of Visual Basic.

Secondly you can use commercially available software, of which probably the most ubiquitous are NinjaTrader, TradeStation and MetaTrader, as well as broker specific packages. I haven't used, and so can't recommend, any particular product. The obvious advantage of these is that they allow you to create an automated strategy without having to write software.

When deciding which to use, the most important criterion is the ability to implement your \textbf{trading rules} within the right \textbf{framework}. \textbf{Back-testing} technology is an added bonus, but make sure it is \textbf{expanding out of sample} or equivalently \textbf{walk forward}. Also it should allow you to fit across multiple \textbf{instruments} and include conservative cost estimates.

Finally you can code up your own trading strategy. In the short term this is far more work, although it's also the most flexible and powerful option. The Interactive Brokers API allows you to use several languages, and there are third-party wrappers available which widen this choice further. For example I use Python to run my strategy via a library which is wrapped around the broker's C++ API. I include some snippets of python code on my website to show you how to implement various parts of the framework.
